%==============================================================================
%  CAPÍTULO XII
%==============================================================================
\chapter{Intersección del derecho concursal con el derecho laboral y el penal económico}
\label{cap:laboral-penal}

\fichaCapitulo{Efectos en los contratos de trabajo y delitos societarios.}{Aplicación del
artículo 163 bis del Código del Trabajo, finiquitos electrónicos y la Ley \Nro 21.595.}

%==============================================================================
\bloqueTeorico
%==============================================================================

\subsection{El derecho laboral y el penal como fronteras del concurso}
\pendiente{Desarrollar.}

\subsection{Fundamento de la preferencia de los créditos laborales}
\pendiente{Desarrollar, en contraste con la persecución penal de la insolvencia
fraudulenta.}

%==============================================================================
\bloqueNormativo
%==============================================================================

\subsection{Área laboral}

\subsubsection{La causal de término del artículo 163 bis}

Causal especial de terminación del contrato de trabajo, que fija como fecha del despido
la de dictación de la resolución de liquidación (\artcdt{163 bis};
\cite{codigo-trabajo}). Sobre el estatuto protector de los trabajadores frente a la
insolvencia del empleador, véase \textcite{doc-proteccion-trabajadores}.

\subsubsection{Preferencia de primera clase del artículo 2472 del Código Civil}

Indemnizaciones por años de servicio, remuneraciones devengadas y feriados
proporcionales, con los topes legales de protección
correspondientes.\verificar{Confirmar los topes vigentes en UF y en ingresos mínimos
mensuales, y su base de cálculo.}

\subsubsection{Finiquitos laborales electrónicos}

Aplicación del instructivo de la \superir{} sobre emisión de finiquitos mediante la
plataforma electrónica de la \dt{} \parencite{instructivo-i6}.\verificar{Número, fecha y vigencia del instructivo.}

\subsection{Área penal}

\subsubsection{Delitos concursales y Ley de Delitos Económicos}

Tipos concursales del \cpen{} y su catalogación por la \Ldeleco{}: ocultación de
información patrimonial relevante y simulación de pasivos, entre
otros \parencite{ley-21595, aj-delitos-economicos}. Para una sistematización del
catálogo, véanse \textcite{doc-delitos-economicos-garrigues} y
\textcite{doc-delitos-economicos-prieto}.\verificar{Confirmar la categoría asignada por la Ley 21.595 y los tipos
comprendidos.}

\subsubsection{Responsabilidad penal de la persona jurídica}

Incidencia de la \Lrpj{} \citeyearpar{ley-20393} y sus modificaciones ante prácticas
delictivas de insolvencia en la administración empresarial.

%==============================================================================
\bloquePractico
%==============================================================================

\subsection{Desvinculación colectiva de los trabajadores}

\begin{flujograma}{Secuencia de desvinculación colectiva tras la resolución de liquidación}{cap12-despidos}
  \node[hito] (sent) {SENTENCIA DE LIQUIDACIÓN DE LA EMPRESA\\[2pt]
    \normalfont Se configura la causal del \artcdt{163 bis}};
  \node[nodo, below=of sent] (comunica)
    {EL LIQUIDADOR COMUNICA EL DESPIDO A LOS TRABAJADORES Y A LA INSPECCIÓN DEL TRABAJO\\[2pt]
     \normalfont Plazo legal de 6 días desde la dictación};
  \node[nodo, below=of comunica] (liquida)
    {LIQUIDACIÓN INDIVIDUAL DE CRÉDITOS Y PREFERENCIAS LABORALES};
  \node[nodo, below=of liquida] (finiquito)
    {EMISIÓN Y SUSCRIPCIÓN DEL FINIQUITO ELECTRÓNICO\\[2pt]
     \normalfont Portal de la \dt{}};
  \node[favorable, below=of finiquito] (pago)
    {PAGO CON PREFERENCIA DE PRIMERA CLASE CON RECURSOS DE LA MASA};

  \draw[flecha] (sent) -- (comunica);
  \draw[flecha] (comunica) -- (liquida);
  \draw[flecha] (liquida) -- (finiquito);
  \draw[flecha] (finiquito) -- (pago);
\end{flujograma}

\begin{alerta}[Plazo de pago del finiquito]
El plazo que la \dt{} aplica ordinariamente para el pago de los finiquitos suscritos en
la plataforma no resulta vinculante en sede concursal: el cobro de las indemnizaciones se
sujeta al estado de realización de los activos de la masa.\verificar{Confirmar que el
instructivo vigente mantiene esta regla.}
\end{alerta}

%==============================================================================
\bloqueJurisprudencial
%==============================================================================

\subsection{Fuero laboral y despido indirecto en sede concursal}

Criterios de los tribunales laborales sobre la vigencia del fuero maternal y sindical de
trabajadores de empresas declaradas en liquidación, y el efecto terminador del
\artcdt{163 bis}.\pendiente{Sistematizar jurisprudencia de unificación.}

\subsection{Responsabilidad de gerentes y directores}

Jurisprudencia que delimita la responsabilidad penal por ocultación o distracción de
activos en los meses previos a la apertura del concurso.\pendiente{Sistematizar.}
